% ---------------------------------------------------------------
\section{Guarantees and Correctness Properties}
\label{sec:guarantees}
% ---------------------------------------------------------------

This section establishes the three foundational correctness guarantees of the
Bailout Protocol: \textbf{SAFETY}, \textbf{LIVENESS}, and \textbf{ACCOUNTABILITY}.
Each guarantee is stated formally, derived from the structural invariants of the
\bxthree{} Framework, and supported by a proof sketch.  Together they constitute
the Bailout Protocol's contribution to the upstream accountability theorem: that
every unresolved exception in a \bxthree{} node reaches a designated human
principal, that contact occurs in finite time, and that the identity of that
principal is permanently attributable.

The guarantees depend on four structural invariants established at node
provisioning and enforced by the \fact{} layer throughout the node's operational
lifetime~\cite{bailoutprotocol2026}:

\begin{itemize}\itemsep=0pt
\item[\textbf{I1}] \textbf{Hierarchy Finiteness.}  The spawn tree depth is bounded
  by $D_{\max}$, a constant established at framework initialization.  Every node
  $n$ has finite depth $d(n) \leq D_{\max}$.

\item[\textbf{I2}] \textbf{Parent-Pointer Immutability.}  The parent pointer
  $\alpha(n)$ is established once at node creation and is read-only for the
  node's entire operational lifetime.  The \fact{} layer's pre-commit hook
  rejects any write attempt to $\alpha(n)$.

\item[\textbf{I3}] \textbf{Human Anchor Immutability.}  The human anchor
  $h(n)$ is designated at node provisioning and is never modified during the
  node's operational lifetime.  No machine actor may redirect the bailout terminus.

\item[\textbf{I4}] \textbf{Forensic Ledger Completeness.}  Every state
  transition in the Bailout Protocol's state machine produces at least one
  Forensic Ledger entry before the transition is committed.  Entries are
  SHA-256 hash-chained and include the identity of every actor that participated
  in the transition.
\end{itemize}

These invariants are not operational conventions.  They are enforced by the
\fact{} layer's deterministic logic below the agent execution level, making them
independent of any machine actor's behavior or intent.

% ---------------------------------------------------------------
\subsection{SAFETY: No Autonomous Collapse}
\label{sec:safety}
% ---------------------------------------------------------------

\textbf{SAFETY} ensures that the system cannot enter a state in which
consequential actions proceed without a human principal in the accountability
chain.  This is the foundational guarantee: if it fails, the other two guarantees
are irrelevant.

\begin{theorem}[Safety --- No Autonomous Collapse]
\label{th:safety}
For any node $n$ in a \bxthree{} spawn tree, the system cannot reach a state
in which autonomous action proceeds without a human principal in the causal chain
of authorization.  Formally,
\[
\forall n \in V,\ \forall s \in Q,\ \text{state}(n) = s \implies
\text{active\_bailout}(n) \lor s \in \{\texttt{IDLE},\ \texttt{RESOLVED}\}.
\]
No terminal state exists that permits autonomous consequential action in the
absence of a human principal.
\end{theorem}

\begin{proof}[Proof sketch]
The proof shows that every pathway to autonomous collapse is structurally blocked
by the conjunction of three mechanisms.

\medskip
\noindent
\textit{Pathway 1: Bypass of the trigger condition.}
The trigger condition TC is evaluated deterministically by the \bounds{} layer
before any autonomous action may proceed.  TC fires when (a) the proposed action
falls outside the node's capability set, (b) the Bounds Engine's projection
forecasts a safety envelope violation above confidence threshold $\tau$, or (c) the
governance policy classifies the action class as human-required.  No machine actor
may suppress, override, or delay the evaluation of TC; it is a pre-commit
condition evaluated by the \fact{} layer before any directive is issued.
If TC fires, the state machine transitions to \texttt{BOUNDED} and all
consequential action stalls until the Bailout Protocol completes or times out.
Autonomous collapse during this interval is blocked.

\medskip
\noindent
\textit{Pathway 2: Suppression of the ESCALATING transition.}
Assume a node reaches \texttt{BAIL\_PENDING} without resolving the trigger.
By Invariant~I2 and Transition T4 (Section~\ref{sec:bp-state-machine}), the
transition $\texttt{BAIL\_PENDING} \xrightarrow{\sigma_{\text{timeout}}}
\texttt{ESCALATING}$ is forced deterministically by the timeout event
$\sigma_{\text{timeout}}$ with no machine-actor discretion.  The \fact{} layer's
pre-commit hook enforces this: any attempt by a machine actor to suppress the
transition or extend the timeout is rejected, the violation is recorded in the
Forensic Ledger, and the Bailout Protocol continues.  No machine actor can
prevent the transition to \texttt{ESCALATING}.

\medskip
\noindent
\textit{Pathway 3: Resolution by a machine actor.}
By the Bypass Property (Definition~\ref{def:bypass}), no machine actor (Plane
P0--P8) can emit the $\sigma_{\text{resolve}}$ event.  The transition
$\texttt{HUMAN\_ACKNOWLEDGED} \xrightarrow{\sigma_{\text{resolve}}}
\texttt{RESOLVED}$ requires the resolver's credential $C_{h(n)}$, which is
provisioned at node creation and stored in the \fact{} layer's credential store.
The pre-commit hook performs this check before committing the state transition:
if the presented credential does not match $C_{h(n)}$, the transition is rejected
and the rejection event is recorded in the Forensic Ledger.  No machine actor
possesses a credential matching $C_{h(n)}$, so no machine actor can emit
$\sigma_{\text{resolve}}$ or otherwise resolve a bailout.

\medskip
\noindent
\textit{Consequence.}
The conjunction of the three pathways establishes that there exists no
sequence of state transitions from \texttt{IDLE} that reaches a state in which
consequential action proceeds without either (i) an active Bailout Protocol in
the \texttt{ESCALATING} or \texttt{HUMAN\_ACKNOWLEDGED} state, or (ii) a
human-authorized \texttt{RESOLVED} state.  $\square$
\end{proof}

The SAFETY guarantee is enforced at the substrate level, not at the agent level.
It does not depend on any machine actor choosing to comply; it depends on the
\fact{} layer's pre-commit hook rejecting any directive that would violate the
safety property.  This makes the guarantee architecturally non-negotiable.

% ---------------------------------------------------------------
\subsection{LIVENESS: Eventual Human Contact}
\label{sec:liveness}
% ---------------------------------------------------------------

\textbf{LIVENESS} ensures that the Bailout Protocol's escalation path from any
node reaches the human anchor $h(n)$ in finite time, or records the failure
attributably if no reachable pathway exists.

\begin{theorem}[Liveness --- Eventual Human Contact]
\label{th:liveness}
For any node $n$ and any trigger $t$ that fires at $n$, the Bailout Protocol's
escalation path $C(n)$ reaches the human anchor $h(n)$ within a bounded number
of state transitions, or transitions to \texttt{RESOLVED} with a
\texttt{PATHWAY\_EXHAUSTED} tag within bounded wall-clock time $T_{\max}$.
Formally,
\[
\forall n \in V,\ \forall t \in T(n),\ \exists\ k \leq d(n) \text{ such that }
\alpha^{(k)}(n) = h(n) \lor \text{state}(n, t) = \texttt{RESOLVED}_{\text{PATHWAY\_EXHAUSTED}}.
\]
\end{theorem}

\begin{proof}[Proof sketch]
The proof follows from the finite termination of the parent-pointer chain and the
deterministic execution of the Escalation Algorithm.

\medskip
\noindent
\textit{Step bound.}  The spawn tree depth is bounded by $D_{\max}$ (Invariant~I1).
The parent-pointer chain from any node $n$ is a strictly ascending sequence in the
depth ordering: for each iteration $i$, $\text{depth}(\alpha^{(i+1)}(n)) >
\text{depth}(\alpha^{(i)}(n))$ because $\alpha^{(i+1)}(n) = \alpha(\alpha^{(i)}(n))$
and $\alpha$ always points to the parent, which has strictly lower depth.  The chain
cannot cycle because self-referential parentage is prevented at node creation by
the \fact{} layer's pre-commit hook.  Therefore the chain from $n$ to the root
terminates in at most $D_{\max}$ steps.

\medskip
\noindent
\textit{Escalation loop termination.}  The escalation algorithm executes one
loop iteration per parent-pointer step.  The loop condition
$\neg\text{IsAnchor}(n_{\text{current}})$ is false precisely when the current
node is $h(n)$, which by the step bound occurs within $D_{\max}$ iterations.
Each iteration writes a \texttt{TRIGGER\_FIRED} or \texttt{ESCALATED} entry to
the Forensic Ledger, providing evidence of participation for every node on the
path.

\medskip
\noindent
\textit{Pathway degradation.}  If an intermediate node $p$ fails to participate
in the escalation, the gap is detected by the Bailout Monitor: the propagation
chain recorded at the anchor will have a missing entry at the expected depth.
The monitor flags this as a pathway violation and may select an alternative
pathway from the Pathway Health Registry, which maintains functional alternative
parent pointers for each node.  The Forensic Ledger records the gap, making the
failure attributable.

\medskip
\noindent
\textit{Terminal failure.}  If all propagation pathways are exhausted (no
functional parent will forward the trigger), the state machine reaches
\texttt{PATHWAY\_EXHAUSTED} and transitions to \texttt{RESOLVED}.  The timeout
event is recorded in the Forensic Ledger with the complete propagation chain
attempted.  This ensures that even under catastrophic failure conditions, the
system records the failure rather than silently suppressing it.

Since the chain terminates in finite time and all failure modes are recorded,
the LIVENESS guarantee holds.  $\square$
\end{proof}

The LIVENESS guarantee is bounded by wall-clock time $T_{\max}$, which is
computed from the provisioned timeout parameters at deployment certification
(Section~\ref{sec:bp-timeouts}).  A deployment that cannot guarantee human
acknowledgment within the required $T_{\max}$ under all failure scenarios must
provision additional or higher-bandwidth pathways, or reduce $T_{\text{prop}}$
and $T_{\text{cap}}$ to bring the worst-case bound within the required limit.
This analysis is a required part of deployment certification for any \bxthree{}
system operating in a regulated or safety-critical environment.

% ---------------------------------------------------------------
\subsection{ACCOUNTABILITY: Human Always Identifiable}
\label{sec:accountability}
% ---------------------------------------------------------------

\textbf{ACCOUNTABILITY} ensures that for every \texttt{RESOLVED} state, the
identity of the resolving human principal is permanently recorded in the Forensic
Ledger and is non-repudiable.

\begin{theorem}[Accountability --- Human Always Identifiable]
\label{th:accountability}
For every trigger $t$ that reaches \texttt{RESOLVED}, the Forensic Ledger
contains a \texttt{human-directive} entry that (i) identifies the resolver by
name and role, (ii) is bound to $t$ via the authorization ledger entry hash,
and (iii) cannot have been produced by any machine actor.  Formally,
\[
\forall t \in T,\ \text{state}(t) = \texttt{RESOLVED} \implies
\exists!\ e \in E.\ e.\text{type} = \texttt{human-directive}
\land e.\text{trigger\_id} = t \land e.\text{resolver} \neq \varnothing
\land e.\text{credential\_match} = \texttt{TRUE}.
\]
\end{theorem}

\begin{proof}[Proof sketch]
The proof follows from three components acting in sequence.

\medskip
\noindent
\textit{Exclusivity of resolution.}  By the Bypass Property
(Definition~\ref{def:bypass}), no machine actor (Plane P0--P8) can emit the
$\sigma_{\text{resolve}}$ event.  The transition
$\texttt{HUMAN\_ACKNOWLEDGED} \xrightarrow{\sigma_{\text{resolve}}}
\texttt{RESOLVED}$ requires the resolver's credential $C_{h(n)}$, which is
provisioned at node creation and stored in the \fact{} layer's credential store.
The pre-commit hook performs this check before committing the state transition:
if the presented credential does not match $C_{h(n)}$, the transition is rejected
and the rejection event is recorded in the Forensic Ledger.  No machine actor
possesses a credential matching $C_{h(n)}$, so no machine actor can emit
$\sigma_{\text{resolve}}$.

\medskip
\noindent
\textit{Forensic Ledger completeness.}  The \fact{} layer enforces
(Invariant~I4) that every state transition produces at least one Forensic Ledger
entry before the transition is committed.  The \texttt{RESOLVED} state
transition therefore necessarily produces a \texttt{human-directive} entry,
which records the resolver's identity (extracted from the authenticated
credential), the directive type, the trigger ID, and the authorization ledger
entry hash that binds the directive to the specific trigger.

\medskip
\noindent
\textit{Non-repudiation.}  The credential used to emit $\sigma_{\text{resolve}}$
is the authoritative credential of $h(n)$, provisioned at node creation and
never shared with any machine actor.  This means the resolution directive is
non-repudiable: $h(n)$ cannot credibly claim that a different actor issued the
directive, because only $h(n)$ possesses the credential required to emit the
event.  The Forensic Ledger's SHA-256 hash chain ensures that the
\texttt{human-directive} entry for trigger $t$ cannot be modified, deleted, or
front-run by a subsequent entry.

Since every \texttt{RESOLVED} state requires a \texttt{human-directive} entry
containing the resolver's identity, and since that identity is credential-bound
and cryptographically chained, the ACCOUNTABILITY guarantee holds.  $\square$
\end{proof}

The ACCOUNTABILITY guarantee extends the upstream accountability theorem beyond
the structural guarantee that human contact is possible, to the operational
guarantee that when human contact occurs, the identity of the human is verifiably
recorded.  This distinction matters in high-stakes deployments: regulators,
customers, and courts require not only that a human was in the loop, but that
the specific human who authorized or rejected an action can be identified with
cryptographic confidence.

% ---------------------------------------------------------------
\subsection{Correctness Triad: Property Interactions}
\label{sec:correctness-triad}
% ---------------------------------------------------------------

The three guarantees interact in essential ways that reinforce the overall
correctness of the system:

\begin{itemize}
\item \textbf{SAFETY $\Rightarrow$ ACCOUNTABILITY.}  By preventing autonomous
  collapse, SAFETY ensures that the attribution chain is always computable ---
  there is always a last known human principal.  A collapsed system cannot
  provide accountability because it no longer has a coherent chain of causation.

\item \textbf{LIVENESS $\Rightarrow$ SAFETY.}  By guaranteeing eventual human
  contact, LIVENESS ensures that the human who can intervene during a SAFETY
  breach is reachable.  Without LIVENESS, SAFETY could be breached indefinitely
  before contact is made.

\item \textbf{ACCOUNTABILITY $\Rightarrow$ LIVENESS enforcement.}  Without
  ACCOUNTABILITY, a malicious or negligent actor could suppress contact signals
  without consequence.  The tamper-evident log creates a non-repudiation
  mechanism that disincentivizes suppression.
\end{itemize}

Together, these three properties form a \textbf{correctness triad} for agentic
systems: the system is safe because it cannot collapse autonomously; it remains
responsive because a human will always eventually be notified or the failure is
recorded; and it remains responsible because every action traces to a person.

\begin{table}[h]
\caption{Guarantees and Correctness Properties: Summary}
\label{tab:guarantees-summary}
\centering
\begin{tabular}{@{}lccc@{}}
\toprule
\textbf{Property} & \textbf{Formal Statement} & \textbf{Mechanism} & \textbf{Critical Assumption} \\
\midrule
SAFETY & $\forall n.\ \text{state}(n) \neq \text{Collapsed}$ &
  Pre-commit hook + Bypass Property & Handoff is non-suppressible \\
LIVENESS & $\exists T_{\max} < \infty.\ \text{contact} \leq T_{\max}$ &
  Progressive escalation + timeout & $\geq 1$ functional pathway exists \\
ACCOUNTABILITY & $\forall t.\ \text{state}(t) = \texttt{RESOLVED} \implies
  \text{resolver}(t) \neq \varnothing$ &
  Credential-bound resolution + hash-chain log & Authentication layer is non-anonymous \\
\bottomrule
\end{tabular}
\end{table}

These properties are not optional enhancements; they are architectural commitments
that must be upheld by the substrate platform and cannot be overridden by any
machine actor at runtime.  Table~\ref{tab:guarantees-summary} summarizes the
three properties, their formal statements, enforcement mechanisms, and critical
assumptions.
